\section{Triangular iteration}\label{ap:iteration}

We finish by organizing the choices in one increment and then iterating them.
The distinction between a large numerical budget and an actual increase in
dimension is essential here.  All constants, polynomial exponents, and
depths of successive polynomial-budget compositions in this section
depend only on $k$.  We use the warm and cold budget conventions of
Definition~\ref{ap:budgets}.

\subsection{Returning to simultaneous residual boxes}

At the root of the ascent, all flags have been removed.  We have the original
input on a root slice, the prepared determining slots, and one fresh absolute
tuple with at most $d_0$ slots.  The determining equations are the exact old
equations after the prescribed changes of integer variables.  In particular,
there is no error to estimate in an argument of a higher determining
polynomial.

\begin{lemma}[Root compression]\label{ap:iteration:compression}
The root return of Section~\ref{ap:extraction} can be converted to the
simultaneous two-box certificate of Theorem~\ref{ap:increment}, with a
further arbitrarily small proportional discount of its target.
The conversion adds no slots.  At weight $i$ its additional logarithmic
width loss and its congruence cost are warm at $i$.  Before common-width
compression, the fresh absolute widths have logarithmic cost polynomial
in $p$ and the dimension bounds.
\end{lemma}

\begin{proof}
Replace each small smooth cutoff by its strict box indicator and each large
smooth cutoff by its enlarged box indicator.  These replacements increase
the positive payoff and decrease the magnitude of the negative payoff.
Reindex the spatial slice and write each determining congruence in the form
$b_i=M_i b_i'+\alpha_i$.  The new residual target is obtained by this exact
substitution and division by $M_i$; its width is divided by the same number.
The logarithm of $M_i$ has the asserted warm cost.

Order the slots of one weight according to their inherited triangular order.
A term involving an earlier slot of this same weight must be linear with
constant coefficient: all variables have positive weights.  By successive
triangular integral changes of the integer slots, reduce those coefficients
to a fixed bounded interval.  Make the same changes exactly in every later
target.  The vector of residuals at this weight now has the form
\[
                    T b-C(u,b_{<i}),
\]
where $T$ is unitriangular, its entries are bounded, and
$\log(2+\|T\|+\|T^{-1}\|)$ is bounded polynomially in the block dimension.
Thus, on putting $\widetilde C=T^{-1}C$, the old residual vector is
$T(b-\widetilde C)$.

Let $P$ be the image under $T^{-1}$ of the old strict residual box, and let
$w_{\min}$ be the least coordinate radius of that old box.  Choose
\[
 0<\rho\le
 \min\left\{\frac1{32},
       \frac{\gamma w_{\min}}{4\max\{1,\|T\|_\infty\}}\right\},
\]
with equality up to fixed factors.  Average centers of simultaneous
$\rho$-boxes in $b-\widetilde C$ over $P+[-\rho,\rho]^m$.
Every old strict point receives full coverage volume $(2\rho)^m$.
Every contribution of a new enlarged box is within
$(2+\gamma)\rho$ of $P$ in maximum norm.  Its image under $T$ is therefore
inside the old enlarged box, since $\gamma\le1$.
The upper enlarged coverage volume at any point is
$(2(1+\gamma)\rho)^m$.  Nonnegativity and a proportional target
discount paying their ratio $(1+\gamma)^m$ select a positive certificate.
This is the same coverage
argument as in Section~\ref{ap:preparation}; it requires neither a lower
bound on the original surplus nor regularity of the sites near a boundary.
It also shows that the selected centers can be kept bounded.

The least of several radii contributes the maximum of their logarithmic
inverse widths.  For old coordinates these logs are the old $Q_i$ plus
warm losses, while for the fresh coordinates they are polynomial in $p$
and dimensions.  Their maximum is therefore at most $Q_i$ plus an allowed
warm loss, padding an empty block with harmless fixed width data.  No new
coordinate is introduced.  We may now reset every value space to its full
ambient space for the next increment.
\end{proof}

\subsection{Structural choices come first}

Section~\ref{ap:extraction} bounds the weights of a fresh absolute patch
by a degree recursion with at most $s$ active steps.  Passive steps do not
increase these weights, and dimensions do not enter the recursion.  Choose
$D$ to dominate its output.  We must also check that this same choice
works in every density round.

Preparation preserves this choice.  A cut at weight $h$ retains a target
of weighted degree at most $h$, introduces a target of degree at most
$h-1$, and replaces the old $h$-slot by an exact polynomial of weighted
degree at most $h$.  Substitution into a higher target consequently preserves
its destination weight.  Slots that were absolute in an earlier round are
simply old determining slots in the next round; those above $s$ are then
passive.  Fix all sampler, family-polynomial, and algebra degree cutoffs
to cover these bounds.  None will subsequently depend on a dimension or
a numerical precision.

\begin{proposition}[Triangular dimension bounds]
\label{ap:iteration:dimensions}
For one round, the increase in the dimension of weight $i$ is bounded by
$d_0$ plus a fixed polynomial in the incoming dimensions strictly above
$i$.  The number of preparation cuts and all structural variable counts
are polynomial in the incoming dimensions and $p$.
For a proposed run of $T=O_k(1+p)$ rounds on an interval, starting without
slots, all these quantities admit polynomial bounds in $2+p$, chosen
before any numerical precision.
\end{proposition}

\begin{proof}
Write $d_i$ for the incoming ambient dimension at weight $i$.  Within one
preparation round set
\[
 H_D=d_D,\qquad H_i=d_i+H_{i+1}^{2}\quad(1\le i<D).
\]
These bound the largest ambient dimensions reached during preparation.
Indeed every cut at $i$ lowers $\dim W_i$ by one.  Its dimension can be
replenished only by incoming full copies from $i+1$, whose dimensions are
already included in $H_i$.  There are therefore at most $H_i$ cuts at $i$.
Each exports a copy of dimension at most $H_i$ to $i-1$, so the total export
is at most $H_i^2$.  This proves the recurrence downwards from $D$.
Restarting a probability test does not reset its value space, so it does
not affect this count.

The fresh absolute tuple adds at most $d_0$ coordinates in total.
Consequently, with $H_{D+1}=0$, the output dimensions satisfy
\begin{equation}\label{ap:iteration:dim-rec}
 d_i^{\rm out}\le d_i+d_0+H_{i+1}^{2}.
\end{equation}
Here and below an expression with $H_{D+1}$ is interpreted as zero.  If a
block is harmlessly relabeled to a still lower weight, replace the last
term by the sum of the corresponding higher exports; the dependence is
still strictly triangular.  The total cut count is at most
$\sum_i H_i$.  Structural path-variable counts are fixed polynomials of
the dimension data, including the root spatial dimension.  Composing
them through the fixed number $D$ of levels preserves polynomial growth.

For completeness, let $P_i(p,d_{>i})$ be nondecreasing polynomials
dominating the additions in \eqref{ap:iteration:dim-rec}, including $d_0$.
Define downward envelopes
\[
 B_D=d_D(0)+T P_D(p),\qquad
 B_i=d_i(0)+T P_i(p,B_{>i}).
\]
An induction first on $i$ downwards and then on the round number gives
$d_i(t)\le B_i$ for $0\le t\le T$.  Starting with $d_i(0)=0$ and
$d_0\le(2+p)^C$, these envelopes are polynomial in $2+p$.
For an arbitrary fixed root spatial dimension $n$, the corresponding
structural forecasts are polynomial in $2+p+n$: the root slices retain
that dimension, and the path-column counts are fixed-degree polynomials
composed through only $D$ layers within each round.  Path dimensions do
not replace the root dimension in the next round.
In the interval application $n=1$.  Choose a polynomial bound $d_*$
dominating the envelopes and all within-round structural counts, and a
polynomial bound $N_c$ for the cut count.  This forecast does not use a
width, rank, scale, or accuracy parameter.
\end{proof}

\subsection{The numerical schedule}

Fix the structural bounds $d_*,N_c$.  Choose the gaps and proportional
discounts so that all preparation selections, the $D$ ascent coverings,
and root compression together consume an arbitrarily small fixed part
of the absolute gain.  For example, $\gamma$ can be a sufficiently small
fixed gain-dependent constant divided by a sufficiently large polynomial
in $2+p+d_*+N_c$.  Divide the permitted preparation discount among the
$N_c$ possible cuts.  Thus $\log(1/\gamma)$ is polynomially bounded before
the remaining choices.  The fixed-depth forward and flag errors may be
chosen as small fixed constants depending on $k$ and the absolute gain.

Let $q_j$ be the incoming width log at level $j$.  We describe a descending
schedule $j=D,D-1,\ldots,1$.  Its record $U_j$ will contain the logarithmic
bounds for every request at level $j$, including its late rank and parent
side requirements.  At the time level $j$ is scheduled, all $U_{>j}$ are
already fixed.  We shall bound each preparation cut's logarithmic width
loss at this level by $A_j$, and the total loss from ascent and root
compression by $A'_j$.  The two required dependencies are
\begin{align*}
 U_j&\le F_j\bigl(p,d_*,N_c,\log(1/\gamma),q_j,U_{>j}\bigr),\\
 N_c A_j+A'_j
    &\le G_j\bigl(p,d_*,N_c,\log(1/\gamma),U_{>j}\bigr),
\end{align*}
for fixed nondecreasing polynomials $F_j,G_j$.  The complete record may
use the current width; the added width loss may not.  We therefore choose
the geometric scales and relative losses first, and only then use $q_j$
to choose the remaining accuracies, rank thresholds, and side cutoffs.
The following four steps verify this order.

\begin{enumerate}
\item \textbf{Normalized geometry and minimum sides.}
Using only $p$, the structural and gap bounds, and $U_{>j}$, choose the
spatial principal scale, the lift ratio $v/w_j$, and the minimum-side
threshold below which an inactive axis is omitted.  These are warm
choices.  The minimum sides must meet the parent-side requirements of
the higher layers, or the absolute-return requirement at the last layer.
They must also permit the full-side separations used when lifting a
rank failure from a strictly higher layer.

In that separation, the path-evaluated slow coefficient and common
denominator bounds, including bases for the later tensor variables,
use only higher-layer data and dimensions.  The separate real and
denominator bounds in Lemmas~\ref{ap:major:refinement}
and~\ref{ap:major:backward} are precisely what ensures this assertion.
Auxiliary subproblems needing cold data at $j$ temporarily set further
pre-short columns to zero; they do not change this minimum-side
requirement.  A rank failure at $j$ itself is lifted only through layers
below $j$.  Thus the present minimum-side choice uses neither the final
rank at $j$, its value-space height, nor its own width loss.

\item \textbf{Preparation width forecast.}
Choose a warm per-cut bound $A_j$, polynomial in the same data and
independent of $q_j$.  In a cut at $h>j$, subdividing the normalized
deviation variables at level $j$ charges the higher failure's coefficient
and denominator bounds, the chosen ratio $v/w_j$, the higher spatial
shrink factors, and the omission cutoff controlling the variation of
omitted rows.  The linear reduction in the incoming copy and the
clearing of its residues charge only higher data in addition to these.
If the lower width must be absolutely smaller than a higher width, the
corresponding relative shrink factor suffices because every old width
is less than one.

A cut at $j$ requires only gap and dimension factors for its own residual
width.  Its large rational section vector is handled in the lower
localizations and the incoming copy, not by conditioning the new
$j$-residual.  The copy can be assigned an arbitrarily small incoming
width without demanding that the formula error be small at that width:
the formula error only needs to be small relative to the gap times the
cut-layer width.  Cuts below $j$ do not shrink level $j$.  We may therefore
budget all preparation widths at this level by
\[
                           q_j+N_c A_j.
\]

\item \textbf{Warm ascent forecast.}
Next forecast the warm modular-product bound, inherited score and slice
bounds, productive probabilities, and macro boundary and cell precisions.
The passive ascent also uses its warm modeling and subslice bounds.
These require only boundedly many polynomial increases in the warm
data, and do not enlarge the preparation bound $A_j$.

At an active step, the micro argument preserves a warm score on the
original warm slice.  The extraction in Section~\ref{ap:extraction}
depends on that score, normalized pivot geometry, the full-law caps
relative to chart mass, and the original slice stride.  It does not
subdivide according to the active micro period, the short-axis lengths,
or the coefficients of the restored polynomial orbit.  The affine-plane
parameterization solves its large congruences identically.  Hence the
relative current-width loss and the returned current-block congruence
cost remain warm.  At passive steps the sharpened major bounds give
the same conclusion.  Include these losses and those of
Lemma~\ref{ap:iteration:compression} in a warm bound $A'_j$.

\item \textbf{Cold-pre and late choices.}
We may now charge $q_j+N_c A_j$, as well as the preceding data.  Forecast
the cold-pre detection, modeling, and major-refinement requests; the
bounded nests of pure-long threshold changes and symbol splittings;
the short-axis and exceptional-product bounds; and every flag-truncation
accuracy required by those requests.  Choose the cutoff for the modulus
cover after the relevant model and period bounds.  Only individual
moduli and witness comparisons of the stated logarithmic size are
required, as in Section~\ref{ap:scalar}.

After all these pre requirements, choose the perturbation, range,
rank, and parent-side bounds.  Choose the rank-stop probability small
enough for the forward calculation using the current prepared width
and higher-layer center masses.  Auxiliary cold computations that
establish a sharpened warm conclusion do not enlarge that conclusion's
output budget.  In particular their own cold costs do not return to
$A_j$ or $A'_j$.

Let $U_j$ dominate this entire schedule and record its rank-failure
requirements for the next lower layer.  Every call to the relative
major refinement at an ancestor is forecast at that ancestor before
its late choices.  The height of $W_j$ after repeated cuts has the
simultaneous stacked-row bound of Proposition~\ref{ap:major:preparation}
in terms of dimensions, cut count, and $\log R_j$; it is not needed in
choosing $R_j$.  Earlier rows are extended
by zeros on incoming copies, and uniform scalar residue changes preserve
the value space.  All rank tests can consequently restart after each
transaction with these uniform schedules.  A nonroot path box pays the
specified later-layer parent-side requirements, not the root slicing
costs of other transactions.
\end{enumerate}

\begin{proposition}[Triangular width bounds]
\label{ap:iteration:widths}
The choices above can be made with fixed polynomial logarithmic bounds.
They give, in one round,
\begin{equation}\label{ap:iteration:width-rec}
 Q_j^{\rm out}\le q_j+P_j(p,d_*,q_{>j}).
\end{equation}
where $P_j$ is a fixed polynomial independent of $q_j$ and of all lower
width logs.  The total root logarithmic slice loss and the sufficient
root log-side cutoff are polynomial in $p,d_*$ and all incoming width
logs.  For the interval run of $O_k(1+p)$ rounds in
Proposition~\ref{ap:iteration:dimensions}, starting without slots and
with trivial width data, these width logs and the sums of root costs
are polynomial in $2+p$.
\end{proposition}

\begin{proof}
Every step in the descending schedule uses a bounded composition of
polynomial logarithmic bounds, giving the two dependencies displayed
above.  Descending substitution through only $D$ levels shows that
$N_c A_j+A'_j$ is polynomial in $p,d_*$ and
$q_{>j}$.  Preparation, ascent, and root compression give
$Q_j^{\rm out}\le q_j+N_c A_j+A'_j$, proving the stated recurrence.

Each transaction's spatial slice loss has polynomial logarithmic cost
in the full recorded budgets.  Their sum over at most polynomially many
transactions has the same form, as does the sufficient root log-side
cutoff.  This statement concerns logarithms: a numerical scale may be
exponential in such a polynomial without violating the estimate.

For this interval run of $T$ rounds starting without slots,
Proposition~\ref{ap:iteration:dimensions} supplies a structural bound
$d_*$ polynomial in $2+p$.  Define downward width envelopes by
\[
 V_D=Q_D(0)+T P_D(p,d_*),\qquad
 V_j=Q_j(0)+T P_j(p,d_*,V_{>j}).
\]
Then $Q_j(t)\le V_j$ for every $t\le T$, by the same descending
induction as for dimensions.  Trivial initial widths, $T=O_k(1+p)$,
and the polynomial bound on $d_*$ make every $V_j$ polynomial in
$2+p$.  The recorded budgets and each
round's root costs are therefore polynomial as well.  Summing the
logarithmic root losses over $T$ rounds preserves this bound.
\end{proof}

\subsection{The increment and the density bound}

\begin{proof}[Completion of the proof of Theorem~\ref{ap:increment}]
Preparation preserves a strict positive certificate with an arbitrarily
small proportional target discount.  Proposition~\ref{ap:scalar:forward}
then supplies productive terminal paths on a
fraction at least a fixed multiple of $a$, at a threshold differing
from $a$ by a sufficiently small fixed proportion.  This use of the
forward mass comparison is important: it does not charge the inverse
of the incoming strict surplus, which may be arbitrarily small.

Apply Lemma~\ref{ap:absolute} on a productive terminal box.  Its sides meet
the dimension-dependent requirement by the minimum-side choices above.
Its fresh rank is bounded by $d_0$ independently of that dimension.
If its hypotheses are impossible, the original certificate cannot
exist.  Otherwise it gives the quantitative positive return at a fixed
multiplicative gain.  The passive and active ascents preserve enough
of that gain, and restore all determining equations exactly.  Apply
Lemma~\ref{ap:iteration:compression} at the root.

Choose the fixed proportional losses sufficiently small that a net
factor $1+\eta_k$, with $\eta_k>0$, remains.  The dimension assertion of
Theorem~\ref{ap:increment} is
Proposition~\ref{ap:iteration:dimensions}; its width assertion and
root-cost assertion are Proposition~\ref{ap:iteration:widths}.
Thus all three required bounds hold simultaneously, with
iteration-independent polynomial exponents.
\end{proof}

\begin{proposition}[Finite-horizon closure]
\label{ap:iteration:closure}
There are constants $C_k>0$ and $A_k\ge1$ such that a progression-free subset of
$[N]$ of density $\alpha>0$ cannot satisfy
\[
                 \log N\ge C_k\bigl(2+\log(1/\alpha)\bigr)^{A_k}.
\]
\end{proposition}

\begin{proof}
Start with no slots, $f=1_A$, and target $a_0=\alpha/2$.
Take one fixed parameter
$p=C(2+\log(1/\alpha))$, with $C$ sufficiently large, and reserve
\[
 T=\left\lceil\frac{\log(2/\alpha)}{\log(1+\eta_k)}\right\rceil+1
       =O_k(1+p)
\]
rounds.  Use this horizon in the advance dimension and width forecasts.
Every temporary application threshold is at least a fixed positive
multiple of $a_0$, so it is at least $e^{-p}$ after enlarging $C$.
The same parameter $p$, fresh-rank bound, and structural forecasts are
therefore valid throughout the run.

Propositions~\ref{ap:iteration:dimensions} and
\ref{ap:iteration:widths} bound the sum of all root logarithmic slice
losses, the maximal remaining log-side requirement, and the endpoint
rounding margins by a fixed polynomial in $2+p$.  Taking the initial
$\log N$ larger than their sum sustains every possible round.  Only
the root is spatially resliced from one round to the next; it remains
an integer-affine slice of the original interval.

If no impossibility occurs, the targets satisfy
$a_t\ge(\alpha/2)(1+\eta_k)^t$, and hence $a_T\ge1$.
But a certificate at any target $a\ge1$ is impossible, since
\[
                         fB^-\le B^+
\]
pointwise.  An absolute-rule impossibility gives the contradiction
sooner.  Enlarging constants to cover the polynomial side requirement
proves the proposition.
\end{proof}

\begin{proof}[Proof of Theorem~\ref{ap:main}]
For $\alpha=r_k(N)/N>0$, the contrapositive of
Proposition~\ref{ap:iteration:closure} gives
\[
 \log N<C_k\bigl(2+\log(1/\alpha)\bigr)^{A_k},
 \qquad
 \alpha<e^2\exp\bigl(-C_k^{-1/A_k}(\log N)^{1/A_k}\bigr).
\]
If $\alpha=0$ the density assertion is immediate.  Renaming the
positive constants gives Equation~\eqref{ap:density-bound}.

We also check the threshold formulation at its endpoints.  Let
$C_k,c_k,\varepsilon_k$ now denote the constants in
Equation~\eqref{ap:density-bound}, and put
$\overline C_k=\max\{1,C_k\}$.  For every $0<\alpha\le1$, the
condition
\[
 \log N\ge
 \left(\frac{\log(2\overline C_k/\alpha)}{c_k}\right)^{1/\varepsilon_k}
\]
implies $r_k(N)\le\alpha N/2<\alpha N$.  Thus a set of density
\emph{at least} $\alpha$ contains the required progression, including
at equality in the density assumption.  Enlarging a constant
$A_k\ge\max\{1,1/\varepsilon_k\}$ makes
$A_k(2+\log(1/\alpha))^{A_k}$ dominate the displayed threshold
for the whole range $0<\alpha\le1$.  Conversely, a threshold of
this form excludes an extremal progression-free set whenever it
holds at $\alpha=r_k(N)/N$, and the preceding inversion gives a
bound of the form Equation~\eqref{ap:density-bound}.

For an integer $r\ge1$, a largest class in an $r$-coloring has
density at least $1/r$.  Substituting $\alpha=1/r$ therefore gives
a monochromatic nonconstant $k$-term progression whenever
\[
 N\ge\left\lceil\exp\bigl(A_k(2+\log r)^{A_k}\bigr)\right\rceil.
\]
The same constant covers $\alpha=1$ and $r=1$.  All constants
depend only on the fixed length $k$.
\end{proof}
